\summarychapter{9}{Laying Her Sceptre Down}{The Conversion of the Empire}

\begin{inbrief}
With Constantine the age of the martyrs ended and two developments began whose effects are still felt:
the Christianizing of the empire and the emperor's interference in the church. The chapter tells how
Diocletian revived the empire and then persecuted the church, how Constantine came to favor it, and
what its new power gained and cost.
\question{Why did a despised and persecuted ``superstition'' rise almost overnight to the spiritual
leadership of the Roman Empire?}
\end{inbrief}

The Arch of Constantine in Rome celebrates his victory over the rival Caesar Maxentius, whose troops
are shown falling from the Milvian Bridge into the Tiber. Eusebius of Caesarea, Constantine's
biographer, compared the scene to Pharaoh's end in the Red Sea, and contemporaries saw in the
emperor's conversion and victory a turning point of history: the opening of what the book calls the Age of the Christian Empire (312--590).

\ssection{Turning Around an Empire}
The story starts with Diocletian. When he took the purple in 284 the empire was collapsing: thirty
emperors in the third century, the Senate ignored, armies and Praetorian Guards making and murdering
rulers. That the collapse was delayed in the West for 175 years, and in the East almost until the
time of Columbus, was largely his work during twenty years of rule (284--305). Because he was the
most savage persecutor he has a bad name among Christians, but Shelley ranks him with the greatest of the emperors. Born to slaves in Dalmatia, he rose through the army, removed his nearest rival
with his own sword and shook the empire back to life, restoring the frontiers on the Rhine and Danube
and recovering Britain and Persia. As a statesman he judged the empire too large for one ruler. He
shared power with three others in four imperial courts, none in Rome (his own at Nicomedia in Asia
Minor): two ``Augusti,'' each with a subordinate ``Caesar'' who could hope to succeed without murder.

\ssection{Christians to the Flames}
For eighteen years Diocletian, a convinced pagan, ignored the Christians. Christian officials filled
his court, his wife Prisca and daughter Valeria were thought to be believers, and large churches stood
in the main cities. Then, two years before the end of his reign and for reasons no one knows, he
purged the army and issued edicts to destroy churches, forbid worship and burn the Scriptures; bishops
were imprisoned, tortured and killed. When he abdicated in 305, forcing his colleague Maximian to do
the same, the persecution went on under Galerius, the new Augustus in the East, whom Christians blamed
for starting it. In the West Constantius Chlorus, who had never pressed it hard, stopped it. The
bloodshed sickened even pagans. Thousands had recanted, but thousands stood firm; some were so eager
to die that the bishop of Carthage refused the title of martyr to those who rushed into it. In 311
Galerius, dying, grudgingly issued an edict of toleration, and Rome's last and worst persecution
ended.

In the struggle for power that followed, Constantine, son of Constantius, crossed the Alps in 312 to
take Italy from Maxentius. Facing a stronger army at the Milvian Bridge outside Rome, he dreamed of a
cross in the sky with the words ``In this sign conquer,'' and his victory on October 28 convinced him
of the power of Christ. Some historians call the conversion purely political: he still plotted and
murdered and kept the pagan title Pontifex Maximus. Shelley thinks his acts say otherwise. From 312
he openly favored Christianity. He exempted its clergy from taxes like pagan priests, abolished
crucifixion, stopped gladiator fights as a punishment, made Sunday a public holiday in 321 and paid
for splendid churches. He brought up his children as Christians and was baptized by Eusebius of
Nicomedia shortly before his death in 337, after which he would wear only his white baptismal robe.
Matthew Arnold pictured the victorious West that ``veil'd her eagles, snapp'd her sword, and laid her
sceptre down.''

\begin{sidebox}{Care for the Dying}
Romans commonly abandoned unwanted children, especially girls. Christians, who saw every person as
made in God's image, rescued them; pagans objected, and laws penalized the rescuers. Constantine ruled
that an abandoned child belonged to its finder, one of many rulings that favored Christians.
\end{sidebox}

\ssection{A New Capital for Christianity}
Constantine also gave the empire a new capital. Its enemies gathered in the east, and its religion
was now eastern. He chose the old town of Byzantium on the Bosporus, a natural crossroads of Asia and
Europe with a perfect harbor, and inaugurated it in 330. As Constantinople it was for most of its
history the capital of the Byzantine, or eastern Roman, Empire; the Turks renamed it Istanbul in 1930.

For the church the change was immense. Outlawed until 312, it was suddenly favored and pampered, and
it re-envisioned its image and mission. Eusebius spoke for most Christians when he presented the
emperor as the ideal Christian ruler and a new age of salvation as begun; many saw providence in a
unified empire prepared for the gospel and thought the conversion of the world near. But there was a
price. Constantine ruled bishops like civil servants and demanded obedience even in purely church
matters. And the masses now streamed in: where the church had consisted of convinced believers who
accepted risk, there came the ambitious, the indifferent and the half-pagan, bringing shallowness,
superstition and the misuse of religion for political ends.

\ssection{The Church Comes to Power}
By 380 rewards for Christians had become penalties for others. In that year the emperor Theodosius
commanded all his peoples to hold the faith of the Trinity and bear the name of Catholic Christians;
the rest, ``demented and insane,'' were heretics whose meeting places were not churches and who would
feel both divine vengeance and his own. He assumed that his will and God's were one, an assumption
built into the Christian empire. Church buildings expressed the new hierarchy: the domed,
mosaic-covered interiors of Byzantine churches, a style borrowed from Persia, displayed God and the
emperor as his representative on earth.

In the West, farther from court, some churchmen dared to challenge the ``demigod.'' In 390 the
people of Thessalonica murdered their governor for refusing to free a popular charioteer he had
jailed. Theodosius, then at Milan, had the crowd at a later race shut in the circus and attacked;
seven thousand died in three hours. Ambrose, bishop of Milan, who saw himself as the emperor's
conscience, wrote to him that he must repent like any other man, since God forgives only the
penitent, and refused him Communion. After staying away for a time Theodosius gave in, laid aside his
imperial robes before the congregation and asked pardon, more than once, until Ambrose gave him the
sacrament on Christmas Day. Ambrose had found the weapon, the threat of excommunication, with which
the Western church would humble princes again and again. In Constantinople no bishop ever went so
far. East and West each claimed the right form of Christianity (one could as well say ``Greek
Catholic'' and ``Roman Orthodox,'' as Bamber Gascoigne notes); their opposite attitudes to the
Christian emperor are a symbol of their diverging destinies.

\begin{keyfacts}
\keypeople{Diocletian; Maximian; Galerius; Constantius Chlorus; Constantine; Maxentius; Eusebius of
Caesarea; Eusebius of Nicomedia; Theodosius; Ambrose of Milan.}
\keydates{%
  \kd{284--305}{Reign of Diocletian; four-man rule; the last great persecution begins about 303}
  \kd{311}{Galerius's edict of toleration}
  \kd{312}{Constantine defeats Maxentius at the Milvian Bridge (October 28)}
  \kd{321}{Sunday made a public holiday}
  \kd{330}{Constantinople inaugurated as the new capital}
  \kd{337}{Baptism and death of Constantine}
  \kd{380}{Theodosius makes Catholic Christianity the religion of the empire}
  \kd{390}{Massacre at Thessalonica; Ambrose forces Theodosius to do penance}}
\keyterms{Augustus and Caesar; edict of toleration; ``In this sign conquer''; Pontifex Maximus;
Constantinople (Byzantium); Christian empire; excommunication.}
\end{keyfacts}
